%! Unit Opener: Functions as Models — Unit 2, Lesson 2.0 — Warm-Up
\documentclass[10pt]{article}
\usepackage{atda-article}
\usepackage{atda-boxes}
\newcommand{\TallMath}[1]{$\displaystyle #1\rule[-1.4em]{0pt}{3.2em}$}
\setlength{\workrowsep}{6pt}   % handwriting room; moves blank and key together

\begin{document}

\pageheader{Unit 2, Lesson 2.0}{Warm-Up --- Unit 2 Diagnostic}
% NAMESTRIP: no \namedateperiod here. The student writes their name once, on the
% cover this component is stapled behind. See references/conventions.md.

\vspace{0.15in}

\begin{notesbox}{Where Are You Starting? --- Five Skills Unit 2 Builds On}
\small
This is a \textbf{diagnostic, not a quiz}. Each item checks one skill this unit is built on. Show
your work; we score it together and you will record the results on your unit roadmap.

\begin{enumerate}[leftmargin=1.6em, itemsep=4pt, topsep=4pt]

\item \textbf{Function notation.} Let $f(x)=-2x+7$. Find $f(-3)$, then find the value of $x$ for
which $f(x)=1$.
\begin{work}
  f(-3) &= -2(-3)+7 \\
        &= 13 \\
  -2x+7 &= 1 \\
    -2x &= -6 \\
      x &= 3
\end{work}

\item \textbf{A value a formula will not take.} For which value of $x$ is
$\dfrac{x+4}{x-5}$ undefined?
\begin{work}
  x-5 &= 0 \\
    x &= 5
\end{work}

\item \textbf{Rate of change.} A table of a car's distance $d$ (miles) after $t$ hours:
\smallskip

\renewcommand{\arraystretch}{1.25}
\begin{tabular}{@{}|c|c|c|c|c|@{}}
\hline
$t$ (hours) & $0$ & $1$ & $2$ & $3$ \\ \hline
$d$ (miles) & $12$ & $20$ & $28$ & $36$ \\ \hline
\end{tabular}
\smallskip

By how much does $d$ change per hour, and what is $d$ when $t=5$?
\begin{work}
  \text{change per hour} &= 20-12 \\
                         &= 8 \text{ miles per hour} \\
                    d(5) &= 12+8(5) \\
                         &= 52 \text{ miles}
\end{work}

\item \textbf{Intercepts of a line.} Find the $x$-intercept and the $y$-intercept of
$3x+2y=12$. Write each as an ordered pair.
\begin{work}
  \text{set } y=0: \quad 3x &= 12 \\
                         x &= 4 \\
  \text{set } x=0: \quad 2y &= 12 \\
                         y &= 6
\end{work}

\item \textbf{Evaluating in context.} A rental car costs \$45 plus \$0.35 per mile driven, so the
cost in dollars for $m$ miles is $C(m)=0.35m+45$. Find $C(80)$, then write one sentence saying
what that number means for the driver.
\begin{work}
  C(80) &= 0.35(80)+45 \\
        &= 28+45 \\
        &= 73
\end{work}
\writeline

\end{enumerate}
\end{notesbox}

\end{document}
